\section{实用决策流程图}

Shikhar Murty 在课程结尾给出了一个非常实用的决策流程图，帮助你在实际项目中选择合适的训练策略：

\begin{figure}[H]
\centering
\includegraphics[width=0.95\textwidth]{slides-images/slide-064.jpg}
\caption{高效训练决策流程图：从混合精度训练开始，逐步尝试 ZeRO Stage 2/3，最后考虑 LoRA\protect\footnotemark}
\end{figure}
\footnotetext{Slides 第65页（最后一页）。}

\begin{importantbox}{高效训练决策流程}
\begin{enumerate}
    \item \textbf{始终使用混合精度训练}。如果 GPU 支持 Ampere 架构（A100/H100/A6000），使用 BF16；否则使用 FP16 + GradScaler
    \item 尝试 \textbf{batch size = 1}。如果能放下：
    \begin{itemize}
        \item 增大 batch size
        \item 始终使用 \textbf{ZeRO Stage 2}（免费优化）
    \end{itemize}
    \item 如果 batch size = 1 也放不下：
    \begin{itemize}
        \item 尝试 \textbf{ZeRO Stage 3 / FSDP}
        \item 尝试 \textbf{Activation Checkpointing}
    \end{itemize}
    \item 如果以上全部失败：使用 \textbf{LoRA}
    \begin{itemize}
        \item 应用于 $W_q$ 和 $W_v$
        \item 秩 $r = 8$，$\alpha = 1$
    \end{itemize}
\end{enumerate}
\end{importantbox}

\begin{warningbox}{此流程假设多 GPU 环境}
上述流程中的 ZeRO / FSDP 假设你有多个 GPU。如果只有单个 GPU，ZeRO 无法帮助你——此时应直接考虑量化（quantization）或 LoRA 等技术。
\end{warningbox}

\subsection{本章小结}

决策流程的核心逻辑是``从简单到复杂''：先用零成本的优化（混合精度、ZeRO Stage 1/2），再考虑有代价的优化（ZeRO Stage 3），最后采用改变训练范式的方法（LoRA）。在每一步都优先选择开销最小的方案。

%% ========== 总结与延伸 ==========

